\def\tituloconcepto{Ruido y la era NISQ}
% Si el capitulo se compila solo (sin main.tex), esta cabecera abre el documento.
% Antes de % Si el capitulo se compila solo (sin main.tex), esta cabecera abre el documento.
% Antes de % Si el capitulo se compila solo (sin main.tex), esta cabecera abre el documento.
% Antes de \input{cabecera} cada capitulo define \tituloconcepto.
\ifdefined\enmain\else
  \documentclass[11pt,a4paper]{article}
  \input{preambulo}
  \begin{document}
  \title{\tituloconcepto}
  \author{\autor}
  \date{\today}
  \maketitle
\fi
 cada capitulo define \tituloconcepto.
\ifdefined\enmain\else
  \documentclass[11pt,a4paper]{article}
  % Preambulo compartido por todos los documentos. Compilar con XeLaTeX/LuaLaTeX
% (o Tectonic). Si se usa Overleaf, elegir el compilador XeLaTeX.
\usepackage[spanish,es-noshorthands]{babel}
\usepackage{amsmath,amssymb}
\usepackage{booktabs,array}
\usepackage[a4paper,margin=2.5cm]{geometry}
\usepackage{xcolor}
\usepackage{listings}
\usepackage[most]{tcolorbox}
\usepackage{quantikz}
\usepackage[hidelinks]{hyperref}

% ---- Autor (editar a gusto) -------------------------------------------------
\newcommand{\autor}{Genesis Silesky}

% ---- Notacion de Dirac y atajos ---------------------------------------------
% quantikz ya define \ket y \bra; DeclareRobustCommand los redefine sin error.
\DeclareRobustCommand{\ket}[1]{\left|#1\right\rangle}
\DeclareRobustCommand{\bra}[1]{\left\langle#1\right|}
\DeclareRobustCommand{\braket}[2]{\left\langle#1\middle|#2\right\rangle}
\newcommand{\C}{\mathbb{C}}
\newcommand{\R}{\mathbb{R}}
\DeclareMathOperator{\Tr}{Tr}

% ---- Cajas de apoyo --------------------------------------------------------
% "concepto": definiciones previas que hacen falta para entender el capitulo.
\newtcolorbox{concepto}[1]{breakable,colback=blue!4,colframe=blue!55!black,
  fonttitle=\bfseries,title={Concepto previo: #1}}
% "nota": aclaraciones y advertencias.
\newtcolorbox{nota}[1]{breakable,colback=orange!6,colframe=orange!70!black,
  fonttitle=\bfseries,title={#1}}

% ---- Estilos de listados ----------------------------------------------------
\lstdefinestyle{py}{language=Python,basicstyle=\ttfamily\footnotesize,
  breaklines=true,postbreak=\mbox{\textcolor{gray}{$\hookrightarrow$}\space},
  numbers=left,numberstyle=\tiny\color{gray},frame=single,
  commentstyle=\color{green!40!black}\itshape,
  keywordstyle=\color{blue!70!black}\bfseries,
  stringstyle=\color{red!60!black},showstringspaces=false,tabsize=4,
  columns=fullflexible,keepspaces=true}
\lstdefinestyle{salida}{basicstyle=\ttfamily\footnotesize,breaklines=true,
  frame=leftline,backgroundcolor=\color{gray!8},columns=fullflexible,
  keepspaces=true}

  \begin{document}
  \title{\tituloconcepto}
  \author{\autor}
  \date{\today}
  \maketitle
\fi
 cada capitulo define \tituloconcepto.
\ifdefined\enmain\else
  \documentclass[11pt,a4paper]{article}
  % Preambulo compartido por todos los documentos. Compilar con XeLaTeX/LuaLaTeX
% (o Tectonic). Si se usa Overleaf, elegir el compilador XeLaTeX.
\usepackage[spanish,es-noshorthands]{babel}
\usepackage{amsmath,amssymb}
\usepackage{booktabs,array}
\usepackage[a4paper,margin=2.5cm]{geometry}
\usepackage{xcolor}
\usepackage{listings}
\usepackage[most]{tcolorbox}
\usepackage{quantikz}
\usepackage[hidelinks]{hyperref}

% ---- Autor (editar a gusto) -------------------------------------------------
\newcommand{\autor}{Genesis Silesky}

% ---- Notacion de Dirac y atajos ---------------------------------------------
% quantikz ya define \ket y \bra; DeclareRobustCommand los redefine sin error.
\DeclareRobustCommand{\ket}[1]{\left|#1\right\rangle}
\DeclareRobustCommand{\bra}[1]{\left\langle#1\right|}
\DeclareRobustCommand{\braket}[2]{\left\langle#1\middle|#2\right\rangle}
\newcommand{\C}{\mathbb{C}}
\newcommand{\R}{\mathbb{R}}
\DeclareMathOperator{\Tr}{Tr}

% ---- Cajas de apoyo --------------------------------------------------------
% "concepto": definiciones previas que hacen falta para entender el capitulo.
\newtcolorbox{concepto}[1]{breakable,colback=blue!4,colframe=blue!55!black,
  fonttitle=\bfseries,title={Concepto previo: #1}}
% "nota": aclaraciones y advertencias.
\newtcolorbox{nota}[1]{breakable,colback=orange!6,colframe=orange!70!black,
  fonttitle=\bfseries,title={#1}}

% ---- Estilos de listados ----------------------------------------------------
\lstdefinestyle{py}{language=Python,basicstyle=\ttfamily\footnotesize,
  breaklines=true,postbreak=\mbox{\textcolor{gray}{$\hookrightarrow$}\space},
  numbers=left,numberstyle=\tiny\color{gray},frame=single,
  commentstyle=\color{green!40!black}\itshape,
  keywordstyle=\color{blue!70!black}\bfseries,
  stringstyle=\color{red!60!black},showstringspaces=false,tabsize=4,
  columns=fullflexible,keepspaces=true}
\lstdefinestyle{salida}{basicstyle=\ttfamily\footnotesize,breaklines=true,
  frame=leftline,backgroundcolor=\color{gray!8},columns=fullflexible,
  keepspaces=true}

  \begin{document}
  \title{\tituloconcepto}
  \author{\autor}
  \date{\today}
  \maketitle
\fi


\section{Ruido y la era NISQ}

\subsection{Objetivo}
Ver cómo los errores de puertas hacen aparecer resultados que en un circuito
ideal serían imposibles ($01$ y $10$ en un estado de Bell).

\subsection{Conceptos previos}
\begin{concepto}{canal de ruido y NISQ}
\begin{itemize}
  \item \textbf{Canal de ruido:} una operación que describe cómo un error
    transforma un estado. El canal \emph{depolarizante} sobre $k$ qubits es
    $\rho\mapsto(1-\lambda)\rho+\lambda\,\tfrac{I}{2^k}$: con probabilidad
    $\lambda$ el estado se reemplaza por el estado máximamente mezclado.
  \item \textbf{NISQ} (\emph{Noisy Intermediate-Scale Quantum}, Preskill 2018):
    dispositivos con decenas a pocos cientos de qubits, con errores, sin
    corrección de errores a escala.
  \item \textbf{Corrección de errores:} codifica un qubit \emph{lógico} en
    muchos qubits \emph{físicos}; el costo (cientos o más por qubit lógico,
    según el código y la tasa de error) es una barrera actual.
\end{itemize}
\end{concepto}

\subsection{Fundamento teórico}
El modelo del código aplica error depolarizante de $1\%$ tras cada $H$ y de
$5\%$ tras cada CNOT. Las puertas de dos qubits tienen más error que las de
uno, lo cual es cualitativamente cierto en hardware real. \textbf{Estos
valores son ilustrativos y no corresponden a ningún dispositivo.}

Con ruido, el estado de salida deja de ser el estado puro de Bell y las
mediciones producen también $01$ y $10$ con probabilidad pequeña. En la
corrida mostrada aparecieron $42$ de $2000$ shots ($2.1\%$) con paridad
incorrecta.

\subsection{Código (completamente comentado)}
\lstinputlisting[style=py]{../08_ruido_nisq.py}

\subsection{Explicación celda por celda}
\begin{description}
  \item[Celda 2] Mismo circuito de Bell; se crea un \texttt{NoiseModel} y se
    asocian errores a las puertas \texttt{h} y \texttt{cx}.
  \item[Celda 3] Se ejecuta el circuito en un simulador ideal y en uno con
    ruido (\texttt{noise\_model=...}) y se comparan los conteos.
\end{description}

\subsection{Salida obtenida}
\lstinputlisting[style=salida]{../salidas/08_ruido_nisq.txt}
Los números exactos cambian en cada ejecución.

\subsection{Interpretación}
El ruido rompe la correlación perfecta. En un equipo real también hay errores
de lectura, decoherencia, etc., que este modelo simple no incluye. Comparar
este histograma con uno obtenido en hardware real (p.\,ej.\ IBM Quantum o las
SpinQ) es la mejor forma de ilustrar la era NISQ.

\subsection{Ejercicios}
\begin{enumerate}
  \item Aumente el error del CNOT a $20\%$ y observe el cambio.
  \item Mida la fidelidad aproximada como la fracción de shots en $00$ o $11$.
  \item Investigue el error de lectura (\texttt{ReadoutError}) y agréguelo.
\end{enumerate}

\ifdefined\enmain\else
  \end{document}
\fi

